\section{Occupational Mapping and Analysis Support}\label{app:mapping}

\subsection{Bridge construction}

The early panel reports Census-2010 occupations and the later panel Census-2018 occupations.  Exposure inputs additionally span SOC vintages.  The versioned bridge first routes 503 early source codes into 568 target codes.  Fifty-six sources have more than one target, with maximum multiplicity seven.  Official shares partition each source record's final weight; separate young and older routing would require validation data that the public adjacent-vintage files do not provide.  The routed early and current stocks reconcile to their sources with relative error below $2\times10^{-16}$.

In the frozen mapping-only lookup, the equal-administrative AIOE variant has full-route coverage for 480 of the 503 pre-2020 source codes.  This 480-code mapping count is not the regression sample: intersecting the rebuilt beta and Webb inputs with positive young preperiod stock yields the 468-occupation primary analysis support below.

The primary joint beta/Webb universe contains 468 occupations and 87.19 percent of stock in the audited candidate universe.  Table~\ref{tab:app_exclusions} gives mutually exclusive exclusion counts.  A machine-readable file lists each included and excluded occupation, its preperiod stocks, scores, and reason.  This file, rather than a hard-coded count in the paper, defines support.

\begin{table}[!htbp]
\centering
\caption{Primary-support exclusions}\label{tab:app_exclusions}
\begin{tabular}{lr}
\toprule
Reason & Occupations \\
\midrule
Missing Rule-A beta only & 41 \\
Missing Webb software only & 22 \\
Missing both exposure inputs & 8 \\
Nonpositive young preperiod stock & 2 \\
\midrule
Total excluded & 73 \\
Primary retained & 468 \\
\bottomrule
\end{tabular}
\end{table}

\subsection{Separating score correction from population expansion}

A literal exact-code merge between SOC-2010 AIOE and post-2018 employment data leaves only 3.33 percent of computer-and-mathematical employment.  The version-aware route reaches 97.7 percent.  The four-row continuous comparison separates the two operations that an aggregate ``crosswalk effect'' would conflate:

The source workbook standardizes AIOE without employment weights across 774 source occupations: its mean is approximately zero and its sample standard deviation is approximately one.  A raw AIOE point is therefore one unit of that source index, not a probability or percentage.  As a common descriptive increment, the source-workbook interquartile difference is 1.879 raw units.  The historical four-row audit below uses a different, explicitly fixed unit: one full-window employment-stock-weighted standard deviation of repaired equal-administrative AIOE on the expanded 495-occupation AIOE-plus-Webb support.  The retained artifact does not store the numerical weighted standard deviation, so I do not manufacture a raw-point conversion.

\begin{table}[!htbp]
\centering
\caption{AIOE taxonomy repair and support expansion}\label{tab:app_crosswalk}
\begin{tabular}{llrr}
\toprule
Score & Support & Occupations & Coefficient per fixed SD \\
\midrule
Literal original & original & 410 & $-0.01885$ \\
Repaired & same original & 410 & $-0.01920$ \\
Repaired & expanded & 495 & $-0.03156$ \\
Repaired & expanded, excluding SOC15 & --- & $-0.02940$ \\
\bottomrule
\end{tabular}
\tabnote{The first row is an invalid vintage merge retained as an implementation audit, not a defensible specification.  Row 1 to 2 isolates score repair on fixed support; row 2 to 3 changes the population; the last row tests whether readmitted computer-and-mathematical occupations alone account for that expansion movement.  Every coefficient uses the fixed weighted-SD unit defined in the text; none is a percentage or a coefficient per raw AIOE point.}
\end{table}

The invalid merge is documented as a failure encountered in this project.  I do not claim that a published paper made the same error without direct verification of its code and occupational universe.

\subsection{Split-route and age-allocation sensitivity}

Within the primary support, split sources contribute 14.93 percent of young and 17.75 percent of older early stock.  Ninety-nine target occupations can receive a split route.  A one-to-one-target restriction and an aggregation to stable Census-2010 categories both retain negative historical-contract estimates, but each changes the target population.  They are mapping sensitivities, not repairs of the primary bridge.

Using the same route proportions by age is transparent but unvalidated.  Accounting ranges that allocate every ambiguous source in all feasible ways place the 2017--2019 young-to-older stock ratio between 0.1031 and 0.1480 in Q1 and between 0.0790 and 0.1175 in Q5.  These are stock-accounting bounds, not bounds on the nonlinear regression coefficient.  An earlier historical-contract tilt involved 6.88 percent of early stock and is retained only as an implementation checkpoint; the current-contract sensitivity follows below.  No adjacent coding vintage is treated as a same-person validation sample.

\subsection{Current-contract route-allocation sensitivity}

The current 113-month contract repeats the exercise on the 468-occupation
support.  Thirty-one split source occupations, accounting for 14.86 percent of
supported early-period routed stock, receive jointly feasible young-versus-older
high/low allocation-odds tilts $K\in\{0.25,0.5,2/3,1,1.5,2,4\}$.  Each grid
point preserves source--age--month mass and official structural route zeros.
Fixed-label rows hold the no-tilt exposure groups constant; rebuilt-treatment
rows recompute them after allocation.  Table~\ref{tab:app_mapping_tilts} keeps
those objects separate.

\begin{table}[!htbp]
\centering
\caption{Current-contract bridge-allocation sensitivity}\label{tab:app_mapping_tilts}
\resizebox{\textwidth}{!}{%
\begin{tabular}{llrr}
\toprule
Model & Exposure labels & Coefficient range & Paired intervals excluding zero \\
\midrule
Pooled & Fixed at $K=1$ & $[-0.1433,-0.1167]$ & 0 of 6 \\
Pooled & Rebuilt at each $K$ & $[-0.1446,-0.1163]$ & 0 of 6 \\
Family by month & Fixed at $K=1$ & $[-0.0748,0.0409]$ & 0 of 6 \\
Family by month & Rebuilt at each $K$ & $[-0.0783,0.0439]$ & 5 of 6 \\
\bottomrule
\end{tabular}}
\tabnote{The grid is $K\in\{0.25,0.5,2/3,1,1.5,2,4\}$, where $K$ is the relative young-versus-older high/low allocation odds for jointly feasible splits.  Ranges summarize coefficients, not confidence intervals.  The last column counts occupation-clustered paired intervals for the six nonbaseline points relative to $K=1$.  All rows are post-outcome mapping sensitivities; the grid is neither an identified set nor a sharp bound.}
\end{table}


All pooled paired occupation-clustered intervals relative to $K=1$ contain
zero under both label rules.  For rebuilt treatment in the family-month model,
however, the coefficient moves from $-0.0217$ at $K=1$ to $0.0439$ at
$K=0.25$; the paired movement is $0.0656$ with interval
$[0.0031,0.1280]$.  Five of the six rebuilt-treatment family-month comparisons
exclude zero, even though rebuilding changes at most two occupations at a grid
point.  Thus the pooled pattern is stable over this declared grid while the
family-conditioned estimand is sensitive to route allocation and resulting
membership.  Null intervals are nondetections, not equivalence.

A separate 25-point grid allows young and older allocation odds to vary jointly
while maintaining the same mass constraints.  The pooled coefficient ranges
from $-0.1599391$ to $-0.0702179$.  No outcome-directed optimization followed
inspection of the grid, but its extrema remain an explored feasible-grid
envelope rather than sharp or global bounds.

Finally, excluding every current occupation that can receive an inbound
one-to-many route leaves 369 occupations.  On this clean-route support, pooled
coefficients are $-0.1642$ ($[-0.2665,-0.0618]$) with fixed labels and
$-0.1468$ ($[-0.2415,-0.0521]$) with rebuilt labels; family-month
coefficients are $-0.0871$ ($[-0.2325,0.0584]$) and $-0.0524$
($[-0.1897,0.0848]$).  On a separate stable-Census-2010 support of 408
occupations, the estimate is $-0.1522$ ($[-0.2438,-0.0606]$).  On the full
468-occupation support restricted to the 77 months from 2020 onward, the
current-label estimates are $-0.1181$ ($[-0.1941,-0.0420]$) pooled and
$-0.0304$ ($[-0.1716,0.1108]$) with family-by-month effects.  These three
exercises change the sample or calendar as well as the labeling rule and are
not a single nested robustness ladder; none validates one coding vintage
against another.

The inherited optional symmetric coding-error sensitivity is not executed:
no authenticated dual-coded CPS sample or validated external misclassification
matrix is available, immediate occupation reversals mix true mobility with
reporting and coding changes, and the route-allocation exercises above are not
relabelled as a coding-error correction.
